\section{Structured SPIN}
\label{sec:structured-spin}

\subsection{From Random Mixing to Structured Mixing}

The examples in Section~\ref{sec:spin-intuition} motivate three changes:
use a fast constituent with a certified spectrum, control how its support
reaches the inner, and amortize random mixing over groups of coordinates.
A sparse message illustrates the design; the proof must also control
larger supports because their multiplicities can dominate the first moment.

First, replace each dense random injection by a constituent with a fast
encoding circuit. Its minimum distance alone is insufficient: nearby
low-weight shells may contain many more words. We require bounds on those
multiplicities. Equal-weight words can still have different coordinate
patterns. Independently permuting each block's coordinates makes
its support uniform conditional on weight, for each fixed message, even
when positions reuse the same constituent.

Second, replace the uniform global interleaver. It can concentrate an
active block near the end of the recursion, leaving too little time to
spread its support. Requiring heavier outer words can reduce that risk
but may increase encoding cost. Instead, transpose the independently
shuffled blocks into regions and shuffle each region independently.
A row of weight $w$ then reaches exactly $w$ distinct regions, with at
most one bit per region. This constrains late activation without
excluding it; the shuffles randomize the remaining alignment.

Third, batch the inner computation. Fixed expansion and feedback maps
combine an input group with a small persistent state, reusing compact
circuits. Each step also applies a sampled transvection, whose action on
a fixed nonzero state mixes a lazy update with uniform refresh.
This Independent-Map Transvection (IMT) inner is analyzed using its actual
state law and the spectra of its two maps, not independent new state bits.
The following construction specifies the maps and setup distributions.

\subsection{Construction at a Glance}

We first define one direct member of the family. Choose a constituent length
$b$, a constituent dimension $k\le b$, a number $L$ of outer-block
positions, an inner step width $t$, and a state dimension $s$, with $t\mid L$.
Set
\begin{equation}
  N:=Lb,
  \qquad
  K:=Lk,
  \qquad
  e:=L/t,
  \qquad
  R:=N/t=be.
  \label{eq:structured-parameters}
\end{equation}
Thus the interleaver will have $b$ regions, each region will contain $L$
coordinates, and the inner will process each region in $e$ steps.

The encoder has the following dataflow:
\begin{equation}
\F_2^{Lk}
  \xrightarrow{\ O\ }
(\F_2^b)^L
  \xrightarrow{\ \Pi_{\mathrm{blk}}\ }
(\F_2^b)^L
  \xrightarrow{\ T\ }
(\F_2^L)^b
  \xrightarrow{\ \Pi_{\mathrm{reg}}\ }
(\F_2^L)^b
  \xrightarrow{\ \mathrm{serialize}\ }
\F_2^N
  \xrightarrow{\ I\ }
\F_2^N.
\label{eq:structured-dataflow}
\end{equation}
Here $O$ applies the constituent encoders, $\Pi_{\mathrm{blk}}$ shuffles
within blocks, $T$ transposes, and $\Pi_{\mathrm{reg}}$ shuffles within
regions. The inner $I$ processes the serialized regions in $t$-bit steps.
The next three subsections define these arrows.

There are two kinds of construction data. The fixed data are the integers
$(L,b,k,t,s)$, two linear maps $A:\F_2^s\to\F_2^t$ and
$C:\F_2^t\to\F_2^s$, a setup law for invertible state maps $M_i:\F_2^s\to\F_2^s$, and a joint setup law for
injective constituent maps $G_j:\F_2^k\to\F_2^b$ for $j\in[L]$.
This joint law must state whether the positions use one fixed constituent, one
shared sampled constituent, or independently sampled constituents. The
distinction changes the expected outer enumerator.

Setup samples the constituent maps from this joint law. Independently, it
samples one uniform coordinate permutation for every outer block, one uniform
permutation for every region, and one state map for
every inner step. All routing permutations and state maps are mutually
independent. After setup, every arrow in
\eqref{eq:structured-dataflow} is fixed. A family selector $\theta_m$ is a
sequence of these parameter choices and setup laws; we suppress $m$ until the
scalable instantiation in Section~\ref{sec:structured-scalable}.

\subsection{Block-Diagonal Outer}

Write a message as $u=(u_1,\ldots,u_L)$ with
$u_j\in\F_2^k$. The outer encoder is the direct sum
\begin{equation}
  O(u)
  :=
  \bigl(G_1(u_1),\ldots,G_L(u_L)\bigr)
  \in(\F_2^b)^L.
  \label{eq:structured-outer}
\end{equation}
Equivalently, a generator matrix for $O$ places the local generator matrices
on the diagonal. The most basic instantiation repeats one constituent
$G:\F_2^k\to\F_2^b$ in every position. The scalable construction below builds
this shared map from a parameterized base constituent and two accumulator
layers. The current certificate uses an extended Golay base constituent, but
the construction interface does not fix that choice or its length.

The outer proof interface is its local weight enumerator. For a realized local
map $G_j$, define
\begin{equation}
  P_{G_j}(z)
  :=
  \sum_{v\in\F_2^k\setminus\{0\}}z^{\wt(G_j(v))}.
  \label{eq:local-constituent-enumerator}
\end{equation}
If one deterministic map $G$ is repeated, the number of messages with exactly
$q$ active blocks and outer weight $h$ is
\begin{equation}
  \binom Lq [z^h]P_G(z)^q.
  \label{eq:repeated-constituent-enumerator}
\end{equation}
For independently sampled, identically distributed local maps, expectation
replaces $P_G$ in \eqref{eq:repeated-constituent-enumerator} by
$\E[P_{G_1}]$. If one sampled map is reused in all positions, the correct
quantity is instead
\begin{equation}
  \binom Lq\E\bigl[[z^h]P_G(z)^q\bigr].
  \label{eq:shared-constituent-enumerator}
\end{equation}
In particular, an expectation over one sampled constituent must not be moved
inside the $q$th power.

The distance proof need not know how the local encoder is implemented. It may
consume an exact enumerator, an expected enumerator under the declared setup,
or a coefficientwise certified envelope. Minimum distance alone is generally
insufficient because nearby low-weight shells can have much larger
multiplicity.

\subsection{Structured Interleaver}

Index the outer coordinates by pairs $(j,r)\in[L]\times[b]$, where $j$ is
the outer-block position. For every $j\in[L]$, setup samples
$\sigma_j\gets S_b$ and defines
$\Pi_{\mathrm{blk}}(j,r):=(j,\sigma_j(r))$.
The deterministic transpose changes the coordinate view by
$T(j,r):=(r,j)\in[b]\times[L]$.
For every region $r\in[b]$, setup samples $\pi_r\gets S_L$ and defines
$\Pi_{\mathrm{reg}}(r,j):=(r,\pi_r(j))$.
The complete interleaver is
\begin{equation}
  \Pi
  :=
  \Pi_{\mathrm{reg}}\circ T\circ\Pi_{\mathrm{blk}}.
  \label{eq:structured-interleaver}
\end{equation}
We serialize the output of \eqref{eq:structured-interleaver} in region-major
order. Thus region $r$ occupies the consecutive coordinates
$(r-1)L+1,\ldots,rL$.

Each outer block contributes exactly one coordinate to each of the $b$
regions. Conditional on the local output weight, $\sigma_j$ makes its active
coordinates a uniform subset of the regions. Conditional on every region
weight, $\pi_r$ makes the support uniform within that region. These two
symmetries remove the original coordinate pattern while preserving the
deterministic spreading supplied by the transpose.

\subsection{Structured Recursive Inner}

The inner processes the serialized interleaver output in $t$-bit steps.
For the IMT recurrences, step inputs, step outputs, and state-space vectors
are columns.
Write that output as $X=(X_0,\ldots,X_{R-1})$, with $X_i\in\F_2^t$ and
$R:=N/t=be$. The state space is $\F_2^s$.
Fix linear maps $A:\F_2^s\to\F_2^t$ and $C:\F_2^t\to\F_2^s$.
Their matrices have dimensions $t\times s$ and $s\times t$, respectively.
Their choice need not satisfy $C=A^\top$ or $CA=0$.
For IMT, setup samples $u_i\gets\F_2^s\setminus\{0\}$ and then
$v_i\gets\{v\in\F_2^s:\langle u_i,v\rangle=0\}$, including zero, and sets
$M_i:=I+u_iv_i^\top$. The pairs are independent over $0\le i<R$.
Since $(u_iv_i^\top)^2=0$, each $M_i$ is invertible.
Starting from $Q_0:=0$, the inner computes
\begin{equation}
  Y_i:=X_i+A(Q_i),
  \qquad
  Q_{i+1}:=M_iQ_i+C(X_i)
  \quad(0\le i<R).
  \label{eq:structured-inner}
\end{equation}
It outputs $I(X):=(Y_0,\ldots,Y_{R-1})$ and discards $Q_R$ without a flush.
State persists across region boundaries. For each fixed $q\ne0$,
\[
  \mathcal L(M_iq)=
  \tfrac12\delta_q+\tfrac12\operatorname{Unif}(\F_2^s\setminus\{0\}).
\]
Here the right side denotes a mixture of probability measures. To see the
identity, $u_i=q$ always fixes $q$; for every other $u_i$, the bit
$\langle v_i,q\rangle$ is fair. Thus any distinct nonzero target has
probability $1/(2(2^s-1))$.

For transposed input groups $U_i\in\F_2^t$, start with $r_R=0\in\F_2^s$
and process steps in reverse, emitting groups $V_i\in\F_2^t$:
\[
 V_i=U_i+C^\top r_{i+1},\qquad
 r_i=A^\top U_i+M_i^\top r_{i+1}.
\]
Both feedback computations consume the raw input, not the emitted output.
The finite constructions below use this same transvection law.

\begin{lemma}[Inner bijectivity]
\label{lem:structured-inner-bijective}
For every fixed state-map schedule and every fixed pair of linear maps
$A,C$, the map $I:(\F_2^t)^R\to(\F_2^t)^R$ defined in
\eqref{eq:structured-inner} is bijective.
\end{lemma}

\begin{proof}
Fix $Y_0,\ldots,Y_{R-1}$. The initial state is known. At step $i$, recover
$X_i=Y_i+A(Q_i)$ and then compute $Q_{i+1}$ from \eqref{eq:structured-inner}. This procedure
recovers exactly one input sequence from every output sequence.
\end{proof}

\subsection{Composition and Rate}

\begin{definition}[Structured SPIN]
\label{def:structured-spin}
For a selector $\theta_m$ and one realization of its setup, the Structured
SPIN encoder is
\begin{equation}
  E_m
  :=
  I_m\circ\Pi_m\circ O_m
  :\F_2^{K_m}\longrightarrow\F_2^{N_m}.
  \label{eq:structured-spin-encoder}
\end{equation}
\end{definition}

Every realized local map is injective. The outer direct sum is therefore
injective. The interleaver is a permutation, and
Lemma~\ref{lem:structured-inner-bijective} makes the inner bijective. Hence
every realized Structured SPIN encoder is injective and has rate
\begin{equation}
  \frac{K_m}{N_m}=\frac{k_m}{b_m}.
  \label{eq:structured-spin-rate}
\end{equation}
The distance analysis can now vary the constituent family without changing
the routing or inner definitions. It must supply a certified outer spectrum
or envelope and a conditional failure envelope for the structured route
classes used in Theorem~\ref{thm:class-first-moment}.

The scalable family below lets the base constituent and $b_m$ vary while
retaining the fixed IMT step and state dimensions. Finite-length
instantiations may tune these choices at one target length. Both cases use
Definition~\ref{def:structured-spin}.
